\documentclass[11pt]{article}
\usepackage[margin=1in]{geometry}
\usepackage{amsmath,amssymb,amsthm,booktabs}
\usepackage[T1]{fontenc}
\usepackage{lmodern,microtype}
\usepackage[colorlinks=true,linkcolor=blue,citecolor=blue,urlcolor=blue]{hyperref}
\newtheorem{theorem}{Theorem}
\newtheorem{lemma}[theorem]{Lemma}
\newtheorem{proposition}[theorem]{Proposition}
\newcommand{\Prob}{\mathbb P}
\newcommand{\C}{\{0,1,2,3\}}
\title{Nonextendibility of symmetric one-dependent\protect\\
four-color laws to a branched ray}
\author{Kaiwen Jin\thanks{Contact: \href{mailto:yui.ui.siki@gmail.com}{yui.ui.siki@gmail.com}.}}
\date{October 13, 2026}
\begin{document}
\maketitle
\begin{abstract}
Let a graph consist of a root, two leaves adjacent to it, and an infinite
ray starting at the root. We prove that this graph admits no color-symmetric,
one-dependent proper four-coloring whose law on every ordered simple path
agrees with a common stationary, reflection-invariant line law. The line law
is arbitrary. An infinite color-avoidance recurrence bounds every specified
one-step extension by half the probability of its prefix. Together with
finite marginal identities, this gives an exact rational contradiction.
The finite calculation is computer-assisted; a complete certificate and a
verifier using only rational arithmetic are supplied.
\end{abstract}

\section{The problem}
One-dependence requires restrictions to vertex sets at distance greater than
one to be independent. It is compatible with proper coloring on the integer
line: Holroyd and Liggett constructed a stationary one-dependent four-coloring
whose law is invariant under color permutations and reflection
\cite[Theorem 4]{HL}. A branch imposes additional consistency conditions on
the same path probabilities.

Liggett and Tang proved a nonextension result for the prescribed
Holroyd--Liggett line law \cite[Proposition 1.4]{LT}. The paragraph following
that proposition distinguishes the question in which the line law is
arbitrary, and Section 6, question 1, asks about that broader possibility.
We address its two-leaf branched-ray case, formulated explicitly as
\href{https://github.com/MColbrook/AIM/blob/c9929805cece044705abaaebe1cebddac1ba11a6/problems/584-four-color-branched-ray.md}{AIM Problem 584}.
All references to \cite{LT} concern its arXiv version 3.

Let $G$ have vertices $o,a,b,v_1,v_2,\ldots$ and edges
$oa,ob,ov_1,v_jv_{j+1}$ for $j\geq1$. A proper four-coloring takes values
in $\C$ and assigns different colors to the endpoints of every edge almost
surely. Graph distance between sets means their minimum vertex distance.

\begin{theorem}\label{thm:main}
There is no pair $(X,Y)$ with the following properties:
\begin{enumerate}
\item $X$ is a random proper four-coloring of $G$, and
$Y=(Y_j)_{j\in\mathbb Z}$ is a random proper four-coloring of the integer line.
\item Each process is one-dependent for arbitrary vertex sets at distance
greater than one, and each law is invariant under every permutation of $\C$.
\item $Y$ is stationary and reflection-invariant.
\item For every ordered simple path $(w_1,\ldots,w_m)$ in $G$,
$(X_{w_1},\ldots,X_{w_m})$ has the law of $(Y_1,\ldots,Y_m)$.
\end{enumerate}
\end{theorem}

No formula for the line law is assumed. Restriction gives the same
nonexistence conclusion if either of the two finite arms is lengthened.
The proof uses the infinite ray, so it does not assert nonexistence on a
fixed finite truncation. It also retains the color and path symmetries in
the theorem.

The infinite ray strengthens the usual finite marginal equations. Avoiding
one fixed color along its continuation gives a recurrence with a double
characteristic root; nonnegativity then bounds any specified extension by
half its prefix probability. Short marginal equations express the unknown
probabilities using two parameters. All products needed at the chosen finite
lengths involve just their square and mixed product. A rational combination
of the resulting identities and bounds is a negative constant. Sections
\ref{sec:cylinders}--\ref{sec:products} derive these facts, and
Section~\ref{sec:certificate} supplies the exact finite contradiction.

\section{Cylinder identities and the infinite tail}\label{sec:cylinders}
Assume such a pair exists. For a color word $w=(w_1,\ldots,w_m)$, write
\[
p(w)=\Prob(Y_1=w_1,\ldots,Y_m=w_m),\qquad p(\varnothing)=1.
\]
For $n\geq0$, write
\[
s(A,B,O,r_1,\ldots,r_n)
=\Prob(X_a=A,X_b=B,X_o=O,X_{v_i}=r_i\ (1\leq i\leq n)).
\]
These are probabilities of specified colors, not masses of color orbits.
Improper words have probability zero. The path condition gives
$s(A,B,O)=p(A,O,B)$.

Color relabeling preserves both types of probability. Stationarity and
reflection also give $p(w)=p(w^{\rm rev})$. To name a color pattern, relabel
its first distinct color as 0, its next distinct color as 1, and so on.
For a line word choose the lexicographically smaller relabeling of the two
orientations. For a branched word only relabel colors: the leaf positions
$a,b$ remain ordered. When several colors in a sum give the same named
pattern, their multiplicities are retained.

Every specified one-site probability is $1/4$, and every proper specified
pair has probability $1/12$. For proper line triples,
\begin{equation}\label{eq:short}
p(i,j,i)=\frac1{48},\qquad p(i,j,k)=\frac1{32}\quad(i\ne k).
\end{equation}
Indeed the endpoints are independent uniforms. For equal endpoints there
are three possible middle colors, and for unequal endpoints there are two;
permutations fixing the endpoints make the allowed choices equiprobable.
The root neighbors $a,b,v_1$ are likewise jointly independent uniforms.
If their specified colors contain $d$ distinct values, then for any allowed
root color
\begin{equation}\label{eq:star-short}
s(A,B,O,r_1)=\frac1{64(4-d)},\qquad d=|\{A,B,r_1\}|.
\end{equation}
Joint independence here follows by applying the arbitrary-set independence
condition successively to a singleton and the remaining set. Permutations
fixing all occupied neighbor colors act transitively on the $4-d$ root choices.

Finite additivity and one-dependence give the following identities. Every
sum is over the four actual colors $c\in\C$; an improper summand is zero.
For line words,
\begin{align}
\sum_c p(uc)&=p(u),\label{eq:line-marginal}\\
\sum_c p(ucv)&=p(u)p(v)\quad(u,v\ne\varnothing).
\label{eq:line-gap}
\end{align}
In \eqref{eq:line-gap}, removing the intervening site separates the two
specified blocks by distance two. Put $r=(r_1,\ldots,r_n)$, $n\geq1$.
The corresponding branched identities are
\begin{align}
\sum_c s(A,B,O,r_1,\ldots,r_{n-1},c)
&=s(A,B,O,r_1,\ldots,r_{n-1}),\label{eq:star-marginal}\\
\sum_c s(A,c,O,r)&=p(A,O,r),&
\sum_c s(c,B,O,r)&=p(B,O,r),\label{eq:leaf}\\
\sum_c s(A,B,c,r)&=\frac1{16}p(r).
\label{eq:root}
\end{align}
The pair of leaves is independent of the entire ray block after removing
the root, which proves \eqref{eq:root}. Finally, for $0\leq i<n$,
\begin{multline}\label{eq:star-gap}
\sum_c s(A,B,O,r_1,\ldots,r_i,c,r_{i+2},\ldots,r_n)\\
=s(A,B,O,r_1,\ldots,r_i)\,p(r_{i+2},\ldots,r_n).
\end{multline}
Deleting $v_{i+1}$ separates the two indicated vertex sets by distance two.
The empty suffix in this formula has probability one.

\begin{lemma}\label{lem:tail}
Let $E$ specify the colors on a finite connected line prefix or a branched
prefix, with an infinite ray continuing beyond its last vertex. For each
color $c$,
\begin{equation}\label{eq:half}
\Prob(E,\text{next vertex has color }c)\leq\frac12\Prob(E).
\end{equation}
\end{lemma}
\begin{proof}
Let $a_n$ be the probability of $E$ followed by $n$ vertices avoiding $c$.
Set $P=\Prob(E)$ and $Q=\Prob(E,\text{next vertex has color }c)$. Then
$a_0=P$ and $a_1=P-Q$. For $n\geq2$, subtract from $a_{n-1}$ the event that
the final vertex has color $c$. On that event its predecessor automatically
avoids $c$, by properness. Dropping the predecessor from the restrictions
leaves the final vertex at distance at least two from all restricted earlier
vertices. Its color is an independent uniform, so
\[
a_n=a_{n-1}-\tfrac14a_{n-2}.
\]
The initial values and induction give
\[
a_n=2^{-n}\bigl[(n+1)P-2nQ\bigr].
\]
If $P-2Q<0$, the bracket is negative for a sufficiently large integer $n$.
This contradicts $a_n\geq0$ and proves the bound, also when $P=0$.
\end{proof}

\section{Short cylinders and their products}\label{sec:products}
Define
\begin{equation}\label{eq:parameters}
\alpha=48p(0,1,0,1),\qquad
\gamma=972s(0,1,2,1,3,1).
\end{equation}
Consider all named proper line words of lengths one through five, and
branched words with one through three ray vertices. There are 18 line and
60 branched probabilities. Apply every identity
\eqref{eq:line-marginal}--\eqref{eq:star-gap} within these ranges, together
with the one-site value and \eqref{eq:parameters}. In every product at least
one factor has the fixed value given above. This is therefore a rational
linear system in the 78 probabilities, with $\alpha,\gamma$ on its right side.

\begin{lemma}[Exact short-cylinder calculation]\label{lem:short-table}
The system just described has rank 78 in its probability columns. Its
solution is affine in $(1,\alpha,\gamma)$, and each line probability is
independent of $\gamma$. In particular,
\begin{align}
p(0,1,0,1,0)&=\frac\alpha{48}-\frac1{288},\label{eq:alpha-lower}\\
s(0,0,1,2,1,3)&=\frac1{1152},&
s(0,0,1,2,1)&=\frac1{256}-\frac\alpha{96},\label{eq:alpha-upper}\\
s(0,0,1,2,1,2)&=\frac\gamma{972}-\frac1{1152},\label{eq:gamma-lower}\\
s(0,1,2,0,2,0)&=\frac1{768}-\frac\gamma{972}.
\label{eq:gamma-upper}
\end{align}
\end{lemma}
\begin{proof}
This is a finite exact linear calculation. In
\path{src/probability_rules.py}, the function \texttt{short\_affine\_table()}
enumerates the specified patterns and constructs the equations from the
displayed identities. Sparse Gaussian elimination normalizes a nonzero
pivot in each probability column, subtracts its rational multiples, and
checks consistency of every redundant row. Back substitution produces all
78 affine expressions. Substitution checks every original equation,
including rows not chosen as pivots. Thus the expressions are the unique
possible values for a hypothetical coloring, rather than one selected
solution of an underdetermined system. The supplied verifier repeats this
calculation using exact fractions; the five entries above and the vanishing
line coefficients of $\gamma$ can also be read from its returned table.
\end{proof}

Nonnegativity in \eqref{eq:alpha-lower}, \eqref{eq:gamma-lower}, and
\eqref{eq:gamma-upper}, and Lemma~\ref{lem:tail} applied to the prefix and
extension in \eqref{eq:alpha-upper}, give
\begin{equation}\label{eq:box}
\frac16\leq\alpha\leq\frac5{24},\qquad
\frac{27}{32}\leq\gamma\leq\frac{81}{64}.
\end{equation}
For the upper bound on $\alpha$, the inequality is explicitly
$1/1152\leq1/512-\alpha/192$.

We now retain all line words of length at most ten and all branched words
with at most eight ray vertices. Set the actual moments
\begin{equation}\label{eq:moments}
t=\alpha^2,\qquad h=\alpha\gamma.
\end{equation}
Every product in \eqref{eq:line-gap} or \eqref{eq:star-gap} becomes affine in
the probability variables and $\alpha,\gamma,t,h$. If either factor is fixed,
this is immediate. Otherwise a line deletion has two factors of length at
least four and total length at most nine. Each thus has length four or five,
and Lemma~\ref{lem:short-table} makes both affine in $\alpha$ alone.
For a ray deletion, a non-fixed branched prefix has at least two ray
vertices and a non-fixed line suffix has at least four vertices. The one
deleted site and the limit of eight ray vertices leave the prefix with at
most three ray vertices and the suffix with at most five vertices. Both
factors are in the short table, and the suffix is independent of $\gamma$.

Explicitly, for rational coefficients with $ge=0$,
\begin{align*}
(c+a\alpha+g\gamma)(d+b\alpha+e\gamma)
={}&cd+(ad+cb)\alpha+(gd+ce)\gamma\\
&+ab\,t+(ae+gb)h.
\end{align*}
This is an identity at the actual moments \eqref{eq:moments}. In particular,
no assertion that arbitrary values of $t,h$ extend to a probability law is
needed.

\section{An exact rational contradiction}\label{sec:certificate}
For a line word of length $m$, specify its odd-position vertices. They are
jointly independent uniforms and their event contains the full word event.
For a branched word with $n$ ray vertices, use both leaves and
$v_1,v_3,v_5,\ldots$. These vertices are also mutually at distance at least
two. Therefore
\begin{equation}\label{eq:independent-upper}
0\leq p(w)\leq4^{-\lceil m/2\rceil},\qquad
0\leq s(A,B,O,r_1,\ldots,r_n)\leq4^{-(2+\lceil n/2\rceil)}.
\end{equation}
Repeated specified colors do not affect this independence argument.
Equations \eqref{eq:box} and \eqref{eq:moments} also give
\begin{equation}\label{eq:moment-upper}
0\leq\alpha\leq\frac5{24},\quad
0\leq\gamma\leq\frac{81}{64},\quad
0\leq t\leq\frac{25}{576},\quad
0\leq h\leq\frac{135}{512}.
\end{equation}

To state the finite calculation, let $F_i$ denote a left side minus right
side of a marginal or deletion identity, the one-site identity, or a
parameter definition. Products are evaluated as above. Let $H_j$ be one
of the following affine expressions:
\begin{itemize}
\item $\Prob(\text{prefix})/2-\Prob(\text{specified one-step extension})$;
\item $t-2c\alpha+c^2$, with $c=5/24$;
\item a probability or one of $\alpha,\gamma,t,h$;
\item its upper bound in \eqref{eq:independent-upper} or
\eqref{eq:moment-upper}, minus that quantity.
\end{itemize}
Every $F_i$ is zero and every $H_j$ is nonnegative for the hypothetical pair.
For the second expression this follows from $(\alpha-c)^2\geq0$.

\begin{proposition}[Exact certificate]\label{prop:certificate}
The supplied file \texttt{certificate/farkas.json} specifies rational
numbers $\lambda_i$ and nonnegative rational numbers $\mu_j$ such that
\begin{equation}\label{eq:certificate}
\sum_{i=1}^{8245}\lambda_iF_i+
\sum_{j=1}^{7244}\mu_jH_j=-D,
\qquad
D=\frac{4663173554073621650668701}{229601458397122150289899520}
>\frac1{50}.
\end{equation}
All probability and moment coefficients on the left cancel exactly.
\end{proposition}
\begin{proof}
The JSON entries identify each equation by its geometry, color word and
deleted position, rather than by a stored matrix row. Each nonnegative
expression identifies its rule and its word or variable. Every multiplier
is a literal rational number. The file has SHA-256
\begin{center}\small\ttfamily
c64ca628b2ea2b324a840827ef7571b9d68f2a8fd4039c7c0baab43581bf4f67
\end{center}
The function \texttt{verify()} in \texttt{src/verify\_certificate.py}
performs the following finite operations:
\begin{enumerate}
\item Derive the short affine table by the rational elimination in
Lemma~\ref{lem:short-table}.
\item Reconstruct each $F_i$ from
\eqref{eq:line-marginal}--\eqref{eq:star-gap}, the one-site value or
\eqref{eq:parameters}; expand its products using \eqref{eq:moments}.
\item Reconstruct each $H_j$ from its stated rule and multiply by its
recorded weight, rejecting a negative $\mu_j$.
\item Add every coefficient and constant with exact fractions. Check that
all 8,298 probability and moment coefficients are zero, that the negative
constant equals the displayed $-D$, and that $D>1/50$.
\end{enumerate}
The nonnegative expressions consist of two tail inequalities, one square,
3,647 lower bounds and 3,594 upper bounds. In particular any small residual
from finding the multipliers is canceled by additional valid bounds in the
exact identity; it is not an arithmetic tolerance. No optimizer or floating
point operation enters this replay. These operations verify
\eqref{eq:certificate} from the supplied finite rational data.
\end{proof}

\begin{proof}[Proof of Theorem~\ref{thm:main}]
The assumed pair supplies its actual finite-cylinder probabilities and
moments. The preceding arguments show that every $F_i$ vanishes and every
$H_j$ is nonnegative at these values. Since every $\mu_j$ is nonnegative,
the left side of \eqref{eq:certificate} is nonnegative. Its right side is
strictly negative by Proposition~\ref{prop:certificate}, a contradiction.
\end{proof}

\paragraph{Reproduction.}
The supplied code and tests accompany the mathematical certificate.
From the supplied directory, run
\begin{quote}\small\ttfamily
python3 src/verify\_certificate.py certificate/farkas.json\\
python3 -m unittest discover -s test -v
\end{quote}
Only the Python standard library is required. The first command repeats the
proof-critical finite calculation; the tests also reject corrupted
certificates and compare short identities with separate probability
controls. The decimal approximation printed for convenience is not used
in a verification decision. These checks accompany the analytic derivations
above; they do not assert independent human or formal verification.

\paragraph{AI-use statement.}
OpenAI Codex was used for literature retrieval, mathematical exploration,
the derivation and checking of the proof, certificate and verifier development,
and manuscript preparation. The finite verification is reproducible with
the supplied rational certificate and code. Independent mathematical review
is requested; no human or Lean certification is claimed.

\begin{thebibliography}{9}
\bibitem{HL}
A.~E. Holroyd and T.~M. Liggett,
\emph{Finitely dependent coloring},
arXiv:1403.2448v4 (2015).
\href{https://arxiv.org/abs/1403.2448v4}{arXiv source}.
\bibitem{LT}
T.~M. Liggett and W. Tang,
\emph{One-dependent colorings of the star graph},
arXiv:1804.06877v3 (2023).
\href{https://arxiv.org/abs/1804.06877v3}{arXiv source}.
\end{thebibliography}
\end{document}
